\section{An unrun transfer study}\label{sec:evaluation}

The proposed study asks whether an object model can describe decisions in other
organisations without concealing cases that do not fit. Its corpus is described
as thirty-two cases, with six node types: requirement, benchmark, bounty,
dependency, call and environment. The intended tests concern encoding residue
and separation of known outcomes. Neither test is reported as passed here.

\begin{figure}
  \centering
  \includegraphics[width=\linewidth]{business-model-grid}
  \caption{Source coverage in the recorded demand-side corpus. Columns locate
  the buyer's problem; colours describe the recorded response. The drawing
  includes cases for which a source locates the problem and lists the others
  below it. These are coverage observations, not results of the proposed
  acceptance experiment. Source:
  \texttt{futon0/analysis/business-models/records/batch-H-demand.edn};
  generator: \texttt{lights.bb}.}
  \label{fig:bm-grid}
\end{figure}

Before running the study, its records need reconciliation. Earlier prose used
both thirty and thirty-two cases, asserted that all cases ended in interest
rather than obligation, and also described a comparison with opposite
terminals. Those statements cannot jointly support a result without a
case-level account of the populations and terminal rule. They are withdrawn as
findings here; reconciliation belongs to the study's preparation.

The proposed discipline is to cite each claim or mark it unverified with a
condition that would settle it. Outcome-blind encoding must be demonstrated by
the collection procedure. Reconstructing what was knowable at the time from
historical accounts is not itself blinding. The project's own contemporaneous
problem statements avoid one source of hindsight, but introduce author
involvement; they are not automatically better evidence overall.

The study may report where the schema fits and where it leaves residue. It
must retain an unmatched case rather than force it into a convenient category.
The stated sample is too small for the proposed predictive claim, and no
result here establishes that acceptance events explain value capture. A
prospective trial would be a separate step, with its criterion and comparison
fixed before the outcome is known.
